% This document was generated by an LLM.

\documentclass{essay}

\title{Logical Foundations of Algebraic Manipulation}
\subtitle{A Guide to Equivalence, Domains, and Reversible Transformations}
\author{}
\date{}

\begin{document}

\maketitle

\section{Introduction}

Elementary algebra is often taught as a collection of symbolic rules.
In reality, every algebraic manipulation is a logical statement. The central
question is not merely whether a step ``looks correct'', but whether it preserves
the mathematical object under consideration.

Throughout mathematics one repeatedly asks:

\begin{itemize}
  \item Has the domain changed?
  \item Has information been lost?
  \item Has information been added?
  \item Is the transformation reversible?
\end{itemize}

These questions unify many topics from elementary algebra to real analysis.

\section{Expressions, Functions, and Domains}

An expression is a symbolic formula together with the values for which it is
defined.

\[ x+1 \]

is defined for every real number.

By contrast,

\[ \frac{x(x+1)}{x} \]

is undefined at $x=0$.

Although

\[ \frac{x(x+1)}{x}=x+1 \]

holds whenever $x\neq0$, the corresponding functions are different because
their domains differ.

\section{Identity versus Equation}

An identity is true for every point in its domain.

\[ (x+1)^2=x^2+2x+1. \]

An equation asks for values satisfying a condition.

\[ x^2+2x+1=9. \]

Confusing identities with equations is a common source of mistakes.

\section{Logical Equivalence and Implication}

The strongest algebraic step is logical equivalence:

\[ P\iff Q. \]

Either statement may replace the other.

A weaker statement is implication:

\[ P\Rightarrow Q. \]

Here solving $Q$ requires checking candidates against $P$.

Example:

\[ a=b\Rightarrow a^2=b^2, \]

but not conversely.

\section{Three Classes of Transformations}

Equivalence transformations preserve the solution set exactly.

Examples include adding the same quantity to both sides or multiplying by a
known nonzero constant.

Implication transformations enlarge the solution set.

Squaring both sides is the classical example.

Conditional transformations are equivalences only under explicit hypotheses.

Dividing by $x$ is equivalent only if $x\neq0$.

\section{Extraneous and Lost Solutions}

Extraneous solutions arise when a transformation enlarges the solution set.

Example:

\[ \sqrt{x}=x-2. \]

Squaring yields

\[ x=(x-2)^2, \]

whose solutions are $1$ and $4$, but only $4$ satisfies the original equation.

Lost solutions arise when information is discarded.

Example:

\[ x(x-1)=0. \]

Dividing by $x$ gives

\[ x-1=0, \]

losing the valid solution $x=0$.

\section{Cancellation}

The cancellation law

\[ ac=bc\Longrightarrow a=b \]

is valid only if $c\neq0$.

This is why

\[ \frac{xc}{c}=x \]

requires the assumption $c\neq0$.

Introducing a denominator therefore introduces a domain restriction that must
remain in force even if the denominator is later cancelled.

\section{Injective and Non-injective Maps}

Many algebraic operations are functions.

If the function is injective, it is reversible.

Examples:

\[ x\mapsto x+5 \]

and

\[ x\mapsto 3x \]

are injective over the real numbers.

However,

\[ x\mapsto x^2 \]

is not injective because

\[ (-2)^2=2^2. \]

This explains why squaring creates extraneous solutions and why square roots
must be handled carefully.

\section{Domains, Codomains, and Partial Functions}

Functions are determined by three pieces of data:

\begin{enumerate}
  \item Domain,
  \item Codomain,
  \item Assignment rule.
\end{enumerate}

Changing any of these changes the function.

Many rational expressions are naturally viewed as partial functions because
they are undefined at certain points.

\section{Why False Proofs Work}

Most famous false proofs exploit a non-reversible transformation.

Typical errors include:

\begin{itemize}
  \item dividing by zero,
  \item cancelling zero,
  \item ignoring domains,
  \item taking square roots incorrectly,
  \item assuming an implication is an equivalence.
\end{itemize}

Each error destroys logical equivalence.

\section{Beyond Algebra}

Exactly the same principles appear throughout mathematics.

In calculus one asks when differentiation and limits commute.

In analysis one asks when sums and integrals may be exchanged.

In linear algebra one asks whether a linear map is invertible.

In abstract algebra one asks in which algebraic structures cancellation holds.

Thus elementary algebra is the first appearance of a general mathematical
theme: understanding which transformations preserve structure.

\section{A Practical Checklist}

Before accepting an algebraic manipulation, ask:

\begin{enumerate}
  \item Has the domain changed?
  \item Is the operation reversible?
  \item Were new assumptions introduced?
  \item Can solutions have been lost?
  \item Can new solutions have been created?
  \item Must every candidate be checked in the original problem?
\end{enumerate}

Developing the habit of asking these questions is more valuable than memorizing
individual rules, because the same reasoning applies from elementary algebra to
graduate mathematics.

\end{document}
